\documentclass[11pt,a4paper]{article}

% Packages
\usepackage[utf8]{inputenc}
\usepackage[spanish, es-tabla]{babel}
\usepackage{caption}
\usepackage{listings}
\usepackage{adjustbox}
\usepackage{enumitem}
\usepackage{hyperref}
\usepackage{boldline}
\usepackage{amsmath}
\usepackage{amssymb, amsmath}
\usepackage[margin=1in]{geometry}
\usepackage{xcolor}
\usepackage{xlop}
\usepackage{soul}
\usepackage[ruled,vlined,linesnumbered]{algorithm2e}
\usepackage{amsthm} %Paquete para la terminología matemática
\usepackage[ruled,vlined,linesnumbered]{algorithm2e}
\usepackage{subfigure}
\usepackage{listings}
\usepackage{color}

\definecolor{dkgreen}{rgb}{0,0.6,0}
\definecolor{gray}{rgb}{0.5,0.5,0.5}
\definecolor{mauve}{rgb}{0.58,0,0.82}

\lstset{frame=tb,
  language=Python,
  aboveskip=3mm,
  belowskip=3mm,
  numbers=left,
  stepnumber=1,
  showstringspaces=false,
  columns=flexible,
  basicstyle={\small\ttfamily},
  numberstyle=\tiny\color{gray},
  keywordstyle=\color{blue},
  commentstyle=\color{dkgreen},
  stringstyle=\color{mauve},
  breaklines=true,
  breakatwhitespace=true,
  tabsize=3
}
%Entorno de la librería matemática (Macros para que no salga en inglés).
\newtheorem{theorem}{Teorema}[section]
\newtheorem{corollary}{Corolario}[theorem]
\theoremstyle{definition}
\newtheorem{definition}{Definición}[section]

% Meta
% Custom
\providecommand{\abs}[1]{\lvert#1\rvert}
\setlength\parindent{0pt}
\definecolor{Light}{gray}{.90}
\newcommand\ddfrac[2]{\frac{\displaystyle #1}{\displaystyle #2}}
\newcommand\tab[1][1cm]{\hspace*{#1}}

\begin{document}
\begin{titlepage}
  \centering
 \includegraphics[width=0.15\textwidth]{./images/exp.jpg}\par\vspace{1cm}
  {\scshape\LARGE Universidad de Granada  \par}
  \vspace{1cm}
  {\scshape Proyecto Final\par}
  \vspace{1.5cm}
  {\huge\bfseries  Clasificación de radiografías torácicas para detección de COVID-19\par}
  \vspace{2cm}
  {\Large\itshape Alberto Luque Infante\\David Villar Martos\par}
  \vfill
  Quinto curso del Doble Grado de Ingeniería Informática y Matemáticas:\par
  Visión por Computador

  \vfill

% Bottom of the page
  {\large \today\par}
\end{titlepage}

\tableofcontents

\section{Descripción del proyecto}

\end{document}



